\documentclass[12pt, a4paper, oneside]{report}

% =========================================================
% PODSTAWOWE PAKIETY
% =========================================================
\usepackage[utf8]{inputenc}
\usepackage[T1]{fontenc}
\usepackage[polish]{babel}
\usepackage{graphicx}
\usepackage{geometry}
\usepackage{setspace}
\usepackage{mathptmx}       % Times New Roman (wymagany)
\usepackage{anyfontsize}    % Do precyzyjnego ustawiania wielkości czcionek
\usepackage{titlesec}       % Do formatowania nagłówków
\usepackage{caption}        % Do podpisów tabel/rysunków
\usepackage[hidelinks]{hyperref} % Linki w spisie treści

% =========================================================
% MARGINESY – WFiIS UŁ
% =========================================================
\geometry{
  a4paper,
  left=30mm,
  right=20mm,
  top=25mm,
  bottom=25mm
}

% =========================================================
% INTERLINIA
% =========================================================
\onehalfspacing

% =========================================================
% START DOKUMENTU
% =========================================================
\begin{document}

% =========================================================
% STRONA TYTUŁOWA – POPRAWIONA WG WZORU 1a
% =========================================================
\begin{titlepage}

% Top logo and faculty name
\begin{center}
\vspace*{0.5cm}
\includegraphics[width=4.69861in,height=1.55278in]{logo.png}
\end{center}

\vspace{1.2cm}

% Name line
\begin{center}
\noindent\makebox[0.85\textwidth]{\dotfill}\\[-0.2em]
{\footnotesize (imi\k{e} i nazwisko)}
\end{center}

\vspace{1.3cm}

% Study info block
\noindent
Kierunek: informatyka\\
Specjalno\'s\'c: informatyka stosowana\\
Specjalizacja: \makebox[0.35\textwidth]{\dotfill}\\
Numer albumu: \makebox[0.35\textwidth]{\dotfill}

\vspace{2.2cm}

% Thesis title line
\begin{center}
\noindent\makebox[0.95\textwidth]{\dotfill}\\[-0.2em]
{\footnotesize (tytu\l{} pracy dyplomowej)}
\end{center}

\vspace{2.6cm}

% Right-side engineering thesis block
\hfill
\begin{minipage}[t]{0.45\textwidth}
\raggedleft
{\bfseries Praca in\.zynierska}\\
wykonana pod kierunkiem\\
\makebox[0.65\textwidth]{\dotfill}\\
w Katedrze \makebox[0.65\textwidth]{\dotfill}\\
WFiIS U\L{}
\end{minipage}

\vfill

% Bottom city and year
\begin{center}
\L{}\'od\'z \makebox[0.25\textwidth]{\dotfill}\\[-0.2em]
{\footnotesize (rok kalendarzowy)}
\end{center}
\end{titlepage}

% =========================================================
% SPISY – NUMERACJA RZYMSKA
% =========================================================
\pagenumbering{roman}

\tableofcontents
\listoffigures
\listoftables

\clearpage

% =========================================================
% TREŚĆ GŁÓWNA – NUMERACJA ARABSKA
% =========================================================
\pagenumbering{arabic}

\chapter{Wstęp}
Tutaj zaczynasz pisać wstęp...

\chapter{Przegląd technologii}
\section{Biblioteka SFML}
Tekst...

\end{document}